\section{Evaluation: Attribution Quality and Optimization Case Studies}
\label{sec:experiments}

We evaluate \sys on two axes: (1) \emph{coverage}---what fraction of kernel runtime is
correctly attributed to \texttt{aten::} operators at each confidence tier; and
(2) \emph{utility}---whether attributed profiles provide sufficient signal to ground and
verify optimization decisions. 

\subsection{Experimental Setup}

\textbf{Hardware.}  All experiments run on a single NVIDIA RTX PRO 6000 Blackwell GPU; full hardware specifications are in \secref{sec:appendix:hwspec}.

\textbf{Workloads.}  \tabref{tab:workloads} describes the three workloads evaluated.  Each
isolates a distinct operator category and hardware challenge.

\begin{table}[h]
\centering
\caption{Case study workloads.}
\label{tab:workloads}
\small
\begin{tabular}{@{}lcp{3.0cm}p{5.2cm}@{}}
\toprule
\textbf{Workload} & \textbf{Batch} & \textbf{Dimensions} & \textbf{Profile Pattern} \\
\midrule
GPT-2 (12-layer)   & 4  & hidden=768, FFN=3072, 12 heads, seq=128 & TC idle (\texttt{tensor\_core\_active\_pct}$\approx$0\%) on all GEMM families \\
SDPA Attention     & 8  & seq=512, dim=512, 8 heads & Low occupancy (16.6\%) + TC idle GEMMs + FP32 attention kernel \\
LSTM Seq.\ Encoder & 32 & seq=128, hidden=512, 2 layers & Structural: Inductor LSTM decomposition into 1,280 scalar SIMT GEMMs \\
\bottomrule
\end{tabular}
\end{table}

\textbf{Measurement protocol.}  Per-operator durations are \nsys CUPTI GPU-domain kernel times.
Each workload runs 2 warm-up iterations followed by 2 measured iterations; reported times are
the sum over measured iterations, and speedup ratios are point estimates without variance
reporting.  At locked clocks, kernel scheduling jitter for these workloads is ${<}$2\% across
repeated runs; point estimates suffice to distinguish the 1.76$\times$--3.63$\times$ speedups
reported here from the 1.0$\times$ null case.  GPU clocks are
locked to the probe-and-cache frequency (\secref{sec:clocklocking}) for the duration of both
the baseline and optimized \nsys captures; both captures lock to the same frequency, so
duration ratios are clock-immune.  \ncu replay is performed in application mode
(\texttt{--replay-mode application}), collecting all 20 metrics in 4--8 passes; \ncu
replay durations are 2--5$\times$ elevated relative to real wall-clock execution and are not
used for speedup reporting.  Layer deduplication is active for models with repeated structure
(\eg GPT-2), profiling unique representatives and propagating metrics to duplicates.

\textbf{Speedup scope.}  All speedup figures in \tabref{tab:speedup} are computed over
\emph{attributed kernel wall time}: the sum of CUPTI GPU-domain durations for kernels that
received an attribution label.  Unattributed kernels (e.g., opaque cuDNN RNN internals in the
LSTM baseline), CPU dispatch overhead, and inter-operator synchronization are excluded.  For
workloads with high unattributed fractions (LSTM: $>$85\% baseline), the reported speedup
reflects the attributed portion only and may differ from end-to-end wall-clock improvement.

\textbf{Baselines.}  Each workload is compared as:
\begin{enumerate*}[label=(\arabic*)]
  \item \emph{Inductor FP32 baseline}: \texttt{torch.compile(backend="inductor")} with
    default settings, FP32 precision;
  \item \emph{Optimized}: workload-specific backend applying BF16 promotion and other
    hardware-evidence-motivated patterns.
\end{enumerate*}
Speedups are reported as \emph{Inductor FP32 baseline} $\div$ \emph{Optimized} in attributed
kernel wall time at locked clocks.

\subsection{Attribution Coverage}
\label{sec:coverage}

Attribution coverage for the three case study workloads is shown in \figref{fig:attribution}.
GPT-2 and SDPA Attention achieve near-complete attribution with the correlation pass active.
LSTM's $>$85\% unattributed fraction is the expected result for that dispatch path; the cause
and corrective strategy are examined in \secref{sec:cs:lstm}.

\begin{figure}[t]
\centering
\begin{tikzpicture}
\begin{axis}[
  xbar stacked,
  width=0.88\columnwidth,
  height=3.6cm,
  bar width=13pt,
  xmin=0, xmax=100,
  xlabel={Kernel runtime attributed (\%)},
  xlabel style={font=\small},
  ytick=data,
  yticklabels={
    LSTM Seq.\ Enc.,
    SDPA Attention,
    GPT-2
  },
  yticklabel style={font=\footnotesize},
  tick label style={font=\footnotesize},
  axis x line=bottom,
  axis y line=left,
  xmajorgrids=true,
  grid style={gray!20},
  legend style={
    at={(0.5,-0.28)}, anchor=north,
    font=\scriptsize, draw=gray!50,
    fill=white, fill opacity=0.9,
    legend columns=3,
  },
]

\addplot[fill=highcol, draw=highcol!80!black] coordinates {
  (4,   0)   % LSTM (opaque cuDNN RNN)
  (35,  1)   % SDPA
  (40,  2)   % GPT-2
};

\addplot[fill=medcol!70, draw=medcol!80!black] coordinates {
  (8,   0)   % LSTM (compiled head fraction)
  (62,  1)   % SDPA
  (57,  2)   % GPT-2
};

\addplot[fill=ucol!80, draw=gray!60] coordinates {
  (88,  0)   % LSTM --- >85% cuDNN RNN opacity
  (3,   1)   % SDPA
  (3,   2)   % GPT-2
};

\legend{HIGH (torch.profiler IOM), MEDIUM (NVTX / Inductor), UNATTRIBUTED}

\end{axis}
\end{tikzpicture}
\caption{\textbf{Attribution coverage for the three case study workloads.}  GPT-2 and SDPA
  Attention achieve $>$95\% attributed runtime with the correlation pass active.  LSTM's
  $>$85\% UNATTRIBUTED fraction correctly reflects opaque cuDNN RNN dispatch
  (see \secref{sec:cs:lstm}).  Inductor fusion enrichment (\secref{sec:heuristic})
  typically adds 5--15 percentage points for compiled workloads.}
\label{fig:attribution}
\end{figure}

\subsection{Profiling Overhead}

\tabref{tab:overhead} characterizes the overhead of the two-stage profiling pipeline, measured
for GPT-2 on RTX PRO 6000 Blackwell.  Layer deduplication (12 transformer blocks $\to$ 1 unique representative)
provides a 12$\times$ reduction in \ncu replay time for this workload.

\begin{table}[h]
\centering
\caption{Two-stage profiling overhead for GPT-2 (RTX PRO 6000 Blackwell, B=4, seq=128, 2 measure
  iterations, locked clocks) with layer deduplication active.  Profiling is an offline, one-time cost.}
\label{tab:overhead}
\small
\begin{tabular}{llll}
\toprule
\textbf{Phase} & \textbf{Wall Time} & \textbf{Overhead} & \textbf{Note} \\
\midrule
Clock probe (first run only)     & ${\sim}$5\,s       & one-time           & Probe-and-cache; reused across all captures \\
\nsys capture (2 iters)          & ${\sim}$10\,s      & $3\times$ baseline & NVTX tracing adds ${\sim}$2$\times$ overhead \\
\ncu replay (unique layers only) & ${\sim}$3--5\,min  & serializes kernels & 12$\times$ speedup vs.\ no dedup (\S\ref{sec:dedup}) \\
Attribution \& aggregation (CPU) & ${<}$5\,s          & negligible         & SQL query + interval tree build \\
\textbf{Total time to profile}   & ${\sim}$4--6\,min  & one-time cost      & Amortized over iterative optimization \\
\bottomrule
\end{tabular}
\end{table}

The profiling cost is commensurate with an offline, iterative optimization workflow: a
single profiling run produces the attributed profile; the optimized backend is implemented
and validated against a second run at matched clock frequencies.  Re-profiling is required
only when the operator set changes between model revisions.

\subsection{Optimization Case Studies}
\label{sec:validation}

The following case studies demonstrate how \texttt{profile.json} grounds optimization
decisions.  Each study identifies the counter symptom read from the attributed profile,
the transformation selected in response, and the before/after counter delta that confirms
bottleneck elimination.  This renders every optimization auditable: the delta report records
\emph{which counter improved}, confirming the diagnosed bottleneck was eliminated rather than
shifted to a different operator.  \tabref{tab:speedup} summarizes attributed kernel wall time
speedups at locked GPU clocks.

\begin{table}[h]
\centering
\caption{Forward-pass attributed kernel wall time and speedup on NVIDIA RTX PRO 6000 Blackwell
 .  All times are from \nsys CUPTI records at locked GPU clocks (not \ncu replay).
  Speedup = Inductor FP32 $\div$ Optimized (attributed kernel time only; see \S\ref{sec:experiments}
  for scope).  Each result is a point estimate from 2 measured iterations.}
\label{tab:speedup}
\small
\setlength{\tabcolsep}{5pt}
\begin{tabular}{lrrr}
\toprule
\textbf{Workload} & \textbf{Inductor FP32 Baseline} & \textbf{Optimized}
                  & \textbf{Speedup} \\
\midrule
GPT-2 (12-layer)        & 7,313\,$\mu$s  & 4,161\,$\mu$s  & \textbf{1.76$\times$} \\
SDPA Attention          & 906\,$\mu$s    & 405\,$\mu$s    & \textbf{2.24$\times$} \\
LSTM Seq.\ Encoder      & 14,531\,$\mu$s & 3,998\,$\mu$s  & \textbf{3.63$\times$} \\
\bottomrule
\end{tabular}
\end{table}

\tabref{tab:hwcounters} reports the key hardware counter deltas that motivated and confirmed
each optimization; each case study below identifies which counter symptom drove the decision.

\begin{table}[h]
\centering
\caption{Key hardware counter deltas for all three case study workloads
  (RTX PRO 6000 Blackwell, \ncu application-mode replay).
  GEMM counters are duration-weighted aggregates (\equref{eq:dwmean}).}
\label{tab:hwcounters}
\small
\setlength{\tabcolsep}{4pt}
\begin{tabular}{llrrl}
\toprule
\textbf{Workload} & \textbf{Counter} & \textbf{Baseline} & \textbf{Optimized}
  & \textbf{Bottleneck Resolved} \\
\midrule
\multicolumn{5}{l}{\textit{GPT-2 (B=4, seq=128)}} \\
& TC active, GEMMs (\%)          & 0.0\%        & 40.5\%       & FP32 SIMT $\to$ BF16 Tensor Core \\
& Registers/thread (GEMMs)       & 210          & 96           & BF16 tiles need fewer registers \\
& Achieved occupancy, GEMMs (\%) & 16.5\%       & ${\sim}$30\% & Register pressure relief \\
& SM throughput, GEMMs (\%)      & ${\sim}$40\% & ${\sim}$65\% & Correct compute pipeline \\
& Attention kernel                & \texttt{fmha\_f32} & \texttt{fmha\_f32} & Stays FP32; SDPA pass removes mask prep \\
\midrule
\multicolumn{5}{l}{\textit{SDPA Attention (B=8, T=512, D=512)}} \\
& TC active, QKV GEMMs (\%)      & 0.0\%  & 19--21\% & FP32 SIMT $\to$ BF16 Tensor Core \\
& Occupancy, QKV GEMMs (\%)      & 16.6\% & 73.9\%   & 3 serial $\to$ 1 wide launch (QKV fusion) \\
& Attention kernel                & \texttt{fmha\_f32} & \texttt{fmha\_bf16} & BF16 promotion activates bf16 fmha \\
\midrule
\multicolumn{5}{l}{\textit{LSTM Seq.\ Encoder (B=32, seq=128, hidden=512)}} \\
& TC active, recurrent GEMMs (\%)  & 0.0\%   & 23.8\% & Scalar SIMT $\to$ cuDNN fused RNN \\
& Occupancy, recurrent GEMMs (\%)  & 9.4\%   & 39.4\% & M=32 $\to$ M=4096 (batch timesteps) \\
& \texttt{splitKreduce} epilogues  & 1{,}280 & 0      & Eliminated by cuDNN routing \\
\bottomrule
\end{tabular}
\end{table}

\subsection{Case Study 1 --- GPT-2: Tensor Core Idle $\to$ BF16 Promotion}
\label{sec:cs:gpt2}

\tabref{tab:transformer} reports per-operator attributed time for GPT-2 small (12 transformer
blocks, hidden=768, 12 heads, FFN=3072, B=4, seq=128) on the RTX PRO 6000 Blackwell.
All times are \nsys-derived GPU kernel wall times at locked clocks (1845\,MHz graphics).
The baseline is dominated by 96 \texttt{cutlass\_80\_simt\_sgemm} kernels; the optimized
backend replaces them with \texttt{cutlass\_80\_tensorop\_bf16} and
\texttt{cutlass\_80\_wmma\_tensorop\_bf16} kernels (\tabref{tab:hwcounters}).

\begin{table}[h]
\centering
\caption{Per-operator-class attributed time for GPT-2 (RTX PRO 6000 Blackwell, B=4, seq=128,
  locked clocks).  Baseline: Inductor FP32 (built-in dedup backend).  Optimized: BF16 dtype
  promotion + SDPA canonicalization + Inductor freezing (\texttt{gpt2\_opt} backend).
  ``Triton casts'' are fp32$\leftrightarrow$bf16 cast kernels introduced by the bf16 promotion
  that the cast-cancellation pass (OPT-4) was unable to eliminate.}
\label{tab:transformer}
\small
\begin{tabular*}{\textwidth}{@{\extracolsep{\fill}}lrrr}
\toprule
\textbf{Operator class} & \textbf{Baseline} & \textbf{Optimized} & \textbf{Speedup} \\
\midrule
GEMM family (\texttt{cutlass\_80\_simt\_sgemm}, 96 kernels) & 6,401\,$\mu$s & 2,695\,$\mu$s & \textbf{2.37$\times$} \\
Attention (\texttt{fmha\_cutlassF\_f32})                    &   485\,$\mu$s &   408\,$\mu$s & 1.19$\times$ \\
LayerNorm                                                   &   149\,$\mu$s &   120\,$\mu$s & 1.24$\times$ \\
Triton casts + elementwise + epilogues                      &   279\,$\mu$s &   938\,$\mu$s & new overhead \\
\midrule
\textbf{Total}                                              & \textbf{7,313\,$\mu$s} & \textbf{4,161\,$\mu$s} & \textbf{1.76$\times$} \\
\bottomrule
\end{tabular*}
\end{table}

The attributed breakdown makes the optimization target unambiguous: 87.5\% of total attributed
runtime (6{,}401 of 7{,}313\,$\mu$s) resides in the GEMM family.  An aggregate
workload-level \texttt{tensor\_core\_active\_pct} reading would confirm that Tensor Cores are
idle, but would not identify whether that idle time is concentrated in GEMMs or distributed
across attention, LayerNorm, and epilogue kernels.  The per-operator breakdown enables the
optimizer to predict, before writing any backend code, that BF16 promotion targeting the GEMM
family will recover a near-proportional fraction of total runtime.

The primary overhead introduced by the optimization is the fp32$\leftrightarrow$bf16 cast
traffic shown in the ``Triton casts'' row: the BF16 promotion pass (\texttt{OPT-1}) inserts
casts around each GEMM node, and the cast-cancellation pass (\texttt{OPT-4}) failed to
eliminate them because \texttt{OPT-1} restores the output dtype directly rather than emitting
per-node inverse cast pairs.  This expanded the Triton elementwise kernel family from 72 to
288 invocations, adding 0.66\,ms and reducing the net attributed-kernel gain from an estimated
${\sim}$2.1$\times$ to the observed 1.76$\times$.

\begin{figure}[t]
\centering
\begin{tikzpicture}
\begin{axis}[
  xbar,
  width=0.88\columnwidth,
  height=4.8cm,
  bar width=14pt,
  xmin=0, xmax=3.2,
  xlabel={Speedup ($\times$, nsys locked-clock wall time)},
  xlabel style={font=\small},
  ytick=data,
  yticklabels={
    {Remaining (LayerNorm + epilogue)},
    {Attention (fmha)},
    {GEMM family}
  },
  yticklabel style={font=\footnotesize},
  tick label style={font=\footnotesize},
  nodes near coords,
  nodes near coords style={font=\scriptsize, anchor=west},
  nodes near coords align={horizontal},
  every node near coord/.append style={xshift=2pt},
  axis x line=bottom,
  axis y line=left,
  xmajorgrids=true,
  grid style={gray!20},
  extra x ticks={1.76},
  extra x tick labels={Total 1.76$\times$},
  extra x tick style={
    grid=major,
    grid style={red!60!black, dashed, thick},
    tick label style={font=\scriptsize, text=red!70!black, rotate=90, anchor=west},
  },
  legend style={
    at={(0.97,0.03)}, anchor=south east,
    font=\scriptsize, draw=gray!50,
    fill=white, fill opacity=0.9,
  },
]

\addplot[fill=orange!70, draw=orange!80!black, bar shift=0pt] coordinates { (2.37, 2) };
\addlegendentry{BF16 GEMM (TC 0\%$\to$40.5\%)}

\addplot[fill=teal!50, draw=teal!70!black, bar shift=0pt] coordinates { (1.19, 1) };
\addlegendentry{BF16 attention}

\addplot[fill=blue!45, draw=blue!60!black, bar shift=0pt] coordinates { (1.22, 0) };
\addlegendentry{BF16 minor gain}

\end{axis}
\end{tikzpicture}
\caption{\textbf{Per-operator-class speedup for GPT-2} (RTX PRO 6000 Blackwell, B=4, seq=128,
  locked-clock nsys wall time).
  \textcolor{orange!80!black}{Orange}: GEMM family (2.37$\times$).
  \textcolor{teal!70!black}{Teal}: attention (1.19$\times$).
  \textcolor{blue!60!black}{Blue}: remaining operators (1.22$\times$).
  Dashed red line: overall \textbf{1.76$\times$}.
  The gap between the GEMM gain and the end-to-end total is explained by new
  fp32$\leftrightarrow$bf16 cast overhead (\tabref{tab:transformer}).}
\label{fig:speedup_bars}
\end{figure}

The attention kernel remains FP32 (\texttt{fmha\_cutlassF\_f32}); the attention speedup
derives from removing the causal-mask construction subgraph (SDPA canonicalization pass),
not a kernel dtype change.  Unlike on A100, the Blackwell attention kernel did not switch to
a lower-register BF16 path---the dominant lever on this workload is exclusively the GEMM
family (\tabref{tab:hwcounters}).

\subsubsection*{Diagnosis and Transformation}

The \texttt{gpt2\_opt} backend applies four optimization passes; the three that produced
measurable counter improvements are described below.  A pass is credited with an improvement
only when its \texttt{status} field is \texttt{APPLIED}, the target counter changes in the
expected direction, and the corresponding operator exhibits reduced runtime.

\textbf{BF16 dtype promotion (OPT-1, +2.37$\times$ on the GEMM family):}
Casting \texttt{mm}/\texttt{addmm} operands to bf16 reroutes all four projection GEMM families
off the FP32 SIMT path and onto the Blackwell BF16 Tensor Core path; hardware counter
confirmation is in \tabref{tab:hwcounters}.

\textbf{SDPA canonicalization (OPT-2):}
Switching \texttt{F.scaled\_dot\_product\_attention} to \texttt{is\_causal=True} removes the
48-launch memory-bound mask-construction subgraph.  This pass runs at the functional level
(\secref{sec:threelevel}) because the high-level SDPA call node is only present before
AOTAutograd decomposition.

\textbf{Inductor freezing (OPT-3):}
\texttt{freezing=True} constant-folds LayerNorm affine weights and projection biases, letting
CUTLASS select fused-epilogue GEMM templates.  The standalone fused-bias Triton family (46
launches, 0.17\,ms) is folded into the GEMM epilogue path in the optimized profile.

\textbf{Cast cancellation (OPT-4, NOT\_APPLIED):}
The cast-cancellation pass detected no per-node inverse cast pairs to cancel because OPT-1
restores the output dtype directly.  The residual fp32$\leftrightarrow$bf16 cast traffic ($+$0.66\,ms,
288 new Triton kernels) is the dominant residual opportunity---revising OPT-1 to maintain
activations in bf16 across consecutive GEMMs would recover most of this overhead, increasing
the net attributed-kernel speedup from the observed 1.76$\times$ toward an estimated
${\sim}$2.1$\times$.

\subsection{Case Study 2 --- SDPA Attention: Low Occupancy and QKV Fusion}
\label{sec:cs:sdpa}

The SDPA workload (multi-head attention, B=8, T=512, D=512, 8 heads) shows a 2.24$\times$
attributed-kernel speedup (\tabref{tab:speedup}).  The baseline concentrates 96.7\% of
attributed time in two operator families: eight \texttt{aten::mm} GEMMs (66.2\%, all TC-idle)
and two SDPA calls (30.5\%, FP32 \texttt{fmha\_cutlassF\_f32\_aligned\_64x64\_sm80}).

A workload-level occupancy reading would reveal low average occupancy but cannot identify
which operator is responsible or why.  The attributed profile isolates occupancy\,=\,16.6\%
to the Q/K/V projection GEMMs specifically, and the FX graph reveals that all three
projections read from a single shared \texttt{LayerNorm} activation---the structural
precondition for QKV fusion.  This causal chain (per-operator counter $\to$ graph-structural
diagnosis $\to$ specific FX pass) is not available from aggregate profiling alone.

Three optimization passes were successfully applied.  \textbf{BF16 dtype promotion} (OPT-1)
rerouted the GEMMs off the FP32 SIMT path and switched the SDPA kernel from
\texttt{fmha\_cutlassF\_f32} to \texttt{fmha\_cutlassF\_bf16}
(${\sim}$\textbf{3.0$\times$} on the attention kernel; counter deltas in \tabref{tab:hwcounters}).
\textbf{QKV fusion} (OPT-2) merged the three Q/K/V projections into a single
\texttt{[512, 1536]} GEMM (three serial 128-block launches $\to$ one 384-block launch),
raising achieved occupancy to 73.9\% (\tabref{tab:hwcounters}).
\textbf{Inductor freezing} (OPT-3) arranged the projection weights in aligned memory
order, enabling the bf16 GEMM autotuner to select aligned templates.

The total attributed-kernel speedup (2.24$\times$) is bounded by newly-introduced auxiliary
kernels---bf16 cast boundaries, the QKV projection split-and-concatenation, and a
pre-transposed frozen-weight conversion---which collectively account for 9.1\% of the
optimized profile runtime.

Unlike the A100 result for FMHA (where register spills were the dominant SDPA bottleneck on
A100), on Blackwell the primary mechanism is Tensor Core activation on the GEMMs and the
bf16 attention kernel, not register pressure reduction.

\subsection{Case Study 3 --- LSTM Sequence Encoder: Structural Miscompilation}
\label{sec:cs:lstm}

The LSTM workload (2-layer LSTM, B=32, seq=128, hidden=512) produces the largest speedup in
the suite (\tabref{tab:speedup}: 3.63$\times$), and uniquely illustrates the framework's
ability to diagnose a structural inefficiency that is undetectable by timing-only profilers.

Under \texttt{torch.compile(backend="inductor")}, \texttt{nn.LSTM} triggers a hard Dynamo
graph break (\texttt{graph\_count = 0}): the operator enters via \texttt{torch.\_VF.lstm}, a
C extension entry point that Dynamo cannot graph-capture.  Inductor therefore decomposes the
LSTM into a per-timestep unrolled loop---1,280 small-matrix FP32 GEMMs (M=32) on the scalar
SIMT path (\texttt{tensor\_core\_active\_pct = 0.0}, achieved occupancy 9.4\%, 1,280 matching
\texttt{cublasLt::splitKreduce\_kernel} epilogues).  Attribution is $>$85\% UNATTRIBUTED because
cuDNN's internal RNN kernels carry no NVTX annotations.

The corrective strategy is structural: given that \texttt{nn.LSTM} unconditionally forces a
Dynamo graph break, the custom backend permits the recurrent region to execute eagerly,
routing it through cuDNN's fused Tensor-Core RNN rather than Inductor's scalar decomposition
path.  cuDNN batches the input-to-hidden projection across all 128 timesteps in one
launch (M=4096 instead of M=32), achieving \textbf{3.73$\times$} on the recurrent GEMMs and
eliminating all 1,280 \texttt{splitKreduce} epilogues, which are absent from the optimized
profile and reduce attributed execution time by 2.26\,ms; hardware counter deltas are in
\tabref{tab:hwcounters}.

This workload illustrates a scenario in which all FX graph passes report
\texttt{NOT\_APPLIED}---because the backend callback is never invoked for the
eagerly-dispatched LSTM region---yet the resulting speedup is the largest observed across
the evaluation suite.  The improvement is attributed to the structural routing decision
(eager dispatch to cuDNN), not to an FX graph transformation; the profile-guided workflow
correctly identified the bottleneck and the optimization follows directly from that
diagnosis.

This distinguishes the attributed pipeline from timing-only profiling.  A kernel-level timer
would expose 1{,}280 slow small-matrix GEMMs and motivate direct kernel optimization---tiling,
BF16 cast, epilogue fusion.  The attribution-tier signal ($>$85\% UNATTRIBUTED from a
cuDNN-opaque path) instead identifies the problem as structural: the LSTM body is not
reachable by any FX-level transformation, and the only valid fix is a dispatch routing
decision.  The attribution tier is the observable that distinguishes these two diagnoses.

